\documentclass[10pt]{book}

\usepackage[T1]{fontenc}
\usepackage{kpfonts}
\usepackage{amsmath}
\usepackage{amssymb}
\usepackage{amsthm}
\usepackage{mathtools}
\usepackage{empheq}
\usepackage{pgfplots}
\usepackage{cleveref}
\usepackage{graphicx}
\usepackage{geometry}
\usepackage{titlesec}
\usepackage{fancyhdr}
\usepackage{setspace}
\usepackage{microtype}
\usepackage{parskip}
\usepackage{tcolorbox}
\usepackage{booktabs}
\usepackage{float}
\usepackage{tikz}

\geometry{a5paper, top=1.5cm, bottom=2cm, inner=1.8cm, outer=1.2cm}
\usetikzlibrary{arrows.meta, positioning}

\titleformat{\chapter}[display]{\normalfont\large\bfseries}{\chaptertitlename\ \thechapter}{10pt}{\large}
\titleformat{\section}{\normalfont\normalsize\bfseries}{\thesection}{1em}{}
\titleformat{\subsection}{\normalfont\small\bfseries}{\thesubsection}{1em}{}

\pgfplotsset{compat=1.18}

\title{Mathematical Theory and Methods}
\author{London Wallace}
\date{Semester 2, 2026}

\renewcommand{\chaptermark}[1]{\markboth{#1}{}}

\pagestyle{fancy}
\fancyhf{}
\fancyfoot[C]{\thepage}
\renewcommand{\headrulewidth}{0pt}
\fancypagestyle{plain}{\fancyhf{}\fancyfoot[C]{\thepage}\renewcommand{\headrulewidth}{0pt}}

\fancypagestyle{plain}{\fancyhf{}\renewcommand{\headrulewidth}{0pt}}

\setstretch{1.15}

\titlespacing{\chapter}{0pt}{-10pt}{10pt}

\newcommand{\newthought}[1]{\vspace{0.5\baselineskip}\noindent{\scshape #1}}

\renewcommand{\bibname}{Bibliography and Suggested Reading}

\begin{document}
\frontmatter
\maketitle
\tableofcontents

\chapter{Preface}
\newthought{These notes} follow, and are meant to cement the course \emph{Algebra}, taught at the
University of Western Australia by Dr. Gordon Royle in Semester 2 of 2026. Where needed, the content of
lectures is supplemented by \emph{Linear Algebra Done Right} by Sheldon Axler.

\section{What is Algebra?}

\newthought{I do not yet know enough to say.}

\section{Table of Symbols}
\begin{center}
  \begin{tabular}{cc}
    \toprule
    Symbol          & Meaning                       \\
    \midrule
    $\forall$       & For every \emph{or} for all   \\
    $\exists$       & There exists                  \\
    $\implies$      & Implies                       \\
    $ \in $         & In \emph{or} is an element of \\
    $ \subset $     & Is a proper subset of         \\
    $ \subseteq $   & Is a subset of                \\
    $ \cup $        & Union with                    \\
    $ \cap $        & Intersection of               \\
    $ \backslash$   & Not \emph{or} without         \\
    $ \varnothing $ & Empty set                     \\
    $ \sim  $       & Relates to                    \\
    $\mapsto$       & Maps to                       \\
    $ > $           & Greater than                  \\
    $ < $           & Less than                     \\
    $ \ge $         & Greater than or equal to      \\
    $ \le $         & Less than or equal to         \\
    $|n|$           & The absolute value of $n$     \\
    $\qedsymbol $   & It has been proven            \\
    
    \bottomrule
  \end{tabular}
\end{center}
\mainmatter
\part{Linear Algebra}

\chapter{Systems of Linear Equations}
\section{Solving Systems of Linear Equations}
\subsection{Augmented Matrices}
\subsection{Elementary Row Operations}
\subsection{Row-Echelon Form}
\section{Gauss-Jordan Elimination}
\section{Reasoning about Systems of Linear Equations}

\chapter{Vector Spaces and Subspaces}
\section{$\mathbb{R}^n$}
\section{Vector Subspaces}
\subsection{Subspace Proofs}
\section{Span}
\section{Linear Independence}
\subsection{The Linear Independence Test}
\section{Bases, Dimension, \& Coordinates}

\chapter{Matrices and Determinants}

\end{document}